\chapter{Réels, intervalles et approximations}
\section*{Exercice 3.1 -- Classer}
Donner le plus petit ensemble parmi $\N,\Z,\D,\Q,\R$ contenant chacun de $-4$, $7/8$, $2/3$, $\sqrt{25}$ et $\sqrt3$. On admet ici l'irrationalité de $\sqrt3$.

\section*{Exercice 3.2 -- Intervalles}
Pour $I=[-3,2[$ et $J=]0,5]$, donner $I\cap J$, $I\cup J$ et $\R\setminus I$.

\section*{Exercice 3.3 -- Distance}
Résoudre $|x+2|\leq3$. Écrire $x\in[4,10]$ à l'aide d'une valeur absolue.

\section*{Exercice 3.4 -- Encadrement}
Encadrer $\sqrt3$ entre deux centièmes consécutifs. Donner son arrondi au centième.

\section*{Exercice 3.5 -- Irrationalité}
En utilisant l'irrationalité de $\sqrt2$, prouver celle de $1+\sqrt2$.

\section*{Exercice 3.6 -- Chiffres significatifs}
Arrondir $0{,}004567$ à trois chiffres significatifs et $98765$ à deux chiffres significatifs. Donner une borne d'erreur absolue dans chaque cas.
